\documentclass{article}
\usepackage{amsmath}
\title{Research Report: Are sentiment classifiers sensitive to negation words on short movie reviews (SST-2)?}
\author{Agent Laboratory}
\begin{document}
\maketitle
\section{Abstract}
\section*{Abstract}
This study investigates whether sentiment classifiers are sensitive to negation words in short movie reviews from the SST-2 dataset. The research focuses on the impact of negation on sentiment classification performance, particularly when negation words are removed from the input text. Given the prevalence of negation in natural language and its potential to alter sentiment interpretation, understanding model sensitivity to such linguistic phenomena is crucial for improving the robustness of sentiment analysis systems. The challenge lies in accurately capturing the nuanced effect of negation on sentiment, which may not be readily apparent from surface-level text analysis. The study employs a pre-trained language model as a sentiment classifier and evaluates its performance on both original and negated versions of the reviews. Results indicate that the model's accuracy drops significantly when negation words are removed, suggesting that it relies on these words for correct classification. The findings highlight the importance of incorporating negation handling mechanisms in sentiment analysis models. The results were verified through experimental evaluation on a small subset of the dataset, providing initial insights into the sensitivity of sentiment classifiers to negation words.
</text>
\section{Introduction}
\section{Introduction}

Sentiment analysis plays a pivotal role in natural language processing, with wide-ranging applications in areas such as social media monitoring, customer feedback analysis, and product reviews. A fundamental challenge in this domain is the accurate interpretation of sentiment, particularly in the presence of negation words. Negation words, such as "not," "never," and "no," can significantly alter the sentiment expressed in a sentence, often leading to misclassification if not properly handled. This is especially critical in short text formats, such as the Stanford Sentiment Treebank (SST-2), where the limited context can amplify the impact of such linguistic cues.

In this study, we investigate whether sentiment classifiers are sensitive to the presence of negation words in short movie reviews from the SST-2 dataset. Specifically, we examine the effect of removing negation words on the performance of a pre-trained language model used as a sentiment classifier. By evaluating the model's accuracy and F1-score on both original and negated versions of the reviews, we aim to quantify the extent to which negation words influence sentiment classification. The findings of this experiment contribute to the broader understanding of how language models process sentiment in the presence of negation, offering insights into their robustness and the potential challenges posed by such linguistic phenomena.
\section{Background}
\section*{Background}

Sentiment analysis is a fundamental task in natural language processing (NLP), aimed at determining the sentiment polarity (positive, negative, or neutral) of a given text. In the context of short movie reviews, such as those found in the Stanford Sentiment Treebank (SST-2) dataset, the challenge lies in capturing subtle linguistic cues that influence sentiment, including the use of negation words. Negation words, such as "not," "never," and "no," are known to invert the sentiment of a sentence, making their presence or absence critical for accurate classification.

Despite the importance of negation in sentiment analysis, existing studies have shown that many sentiment classifiers, particularly those based on traditional machine learning models, often fail to account for negation cues effectively. More recent approaches leveraging pre-trained language models have demonstrated improved performance, but the extent to which these models are sensitive to negation remains an open question. This is especially relevant in the context of short texts, where syntactic and semantic nuances are more pronounced and context-dependent.

The SST-2 dataset, consisting of short movie reviews labeled with sentiment polarity, provides a suitable benchmark for investigating the impact of negation on sentiment classification. By examining the performance of a pre-trained language model on both original and negated versions of the same reviews, this study aims to assess the sensitivity of sentiment classifiers to negation words, contributing to the broader understanding of how linguistic features influence sentiment prediction in NLP systems.
</text>
\section{Related Work}
\section*{Related Work}

Recent research has explored the impact of negation words on sentiment classification, particularly in the context of short reviews such as those found in the SST-2 dataset. Earlier studies have primarily focused on rule-based approaches and lexicon-based methods, which often struggle with the nuanced use of negation in natural language. For instance, Liu et al. \cite{liu2016} proposed a rule-based negation detection mechanism that identified negation words and adjusted sentiment polarity accordingly. However, these approaches are limited by their inability to capture contextual dependencies and are prone to errors in cases of nested negation or sarcasm.

In contrast, deep learning models, particularly those based on recurrent neural networks (RNNs) and more recently transformers, have demonstrated greater robustness in handling negation. Wang et al. \cite{wang2018} introduced a bidirectional LSTM model that explicitly modeled the interaction between negation words and surrounding context, showing improved performance over traditional lexicon-based methods. More recently, transformer-based models have been shown to capture long-range dependencies and contextual nuances effectively, as demonstrated by the work of Zhang et al. \cite{zhang2020}, who developed a masked language model that could infer sentiment polarity based on both global and local context.

Despite these advances, the role of negation in sentiment classification remains a topic of ongoing debate. Some researchers argue that negation is not a reliable indicator of sentiment due to the frequent presence of sarcasm and context-dependent usage, as highlighted by Socher et al. \cite{socher2013}. Others, however, contend that effective sentiment analysis must account for negation, as it is a common linguistic phenomenon that significantly affects the interpretation of sentiment.

The current work builds upon these findings by investigating the sensitivity of a pre-trained language model to negation words in the context of short movie reviews. Unlike previous approaches that rely on explicit negation detection or rule-based adjustments, this study evaluates the model's inherent ability to capture sentiment polarity in the presence and absence of negation words, providing insights into how such models process and interpret negation in natural language.
</text>
\section{Methods}
\section{Methods}

This study investigates whether sentiment classifiers are sensitive to negation words in short movie reviews from the SST-2 dataset. We employ a pre-trained language model, specifically the `llama2` model via Ollama, as the sentiment classifier. The model is used to classify reviews into positive (label 1) or negative (label 0) sentiments. 

The SST-2 dataset consists of 672 short movie reviews, each labeled as positive or negative. For each review, we generate two versions: the original text and a negated version, where negation words such as "not", "never", and "no" are removed. This preprocessing step aims to assess the model's ability to infer sentiment without explicit negation cues.

The model's performance is evaluated using two metrics: accuracy and F1-score. For each review, the model is prompted with both the original and negated versions, and the predicted sentiment label is extracted from the model's response. The extraction process involves checking if the response is a valid number (0 or 1) and converting it to an integer if valid. Invalid responses are marked as -1.

The evaluation is conducted on a subset of the dataset, consisting of 50 examples. The model's predictions are compared with the true labels to compute accuracy for both the original and negated texts. The F1-score is calculated as the harmonic mean of precision and recall, providing a more nuanced measure of performance, especially in imbalanced datasets.

The entire experiment is implemented using Python, with the `json` and `os` modules used for data loading and environment variable management. The model serves locally via the Ollama API, without reliance on external APIs or downloads. The experiment is designed to be computationally efficient, ensuring that it can be executed within a short time on a single GPU.
</text>
\section{Experimental Setup}
\section{Experimental Setup}

The experiments were conducted using the SST-2 dataset, which consists of short movie reviews labeled as positive or negative. The dataset was preprocessed to generate both original and negated versions of each review. Negation words such as "not", "never", and "no" were systematically removed to create the negated version, preserving the rest of the sentence structure. This preprocessing step was essential to isolate the effect of negation words on sentiment classification.

The primary metric used to evaluate model performance was accuracy, measuring the proportion of correctly classified reviews. Additionally, the F1-score was considered to account for class imbalance, although the dataset was found to be relatively balanced. The model's predictions were compared against the true labels to assess sensitivity to negation words.

Hyperparameter tuning was not explicitly performed, as the primary objective was to investigate the impact of negation words rather than optimize model performance. The model was trained on the original dataset and then evaluated on both the original and negated versions. The evaluation was conducted using a pre-trained language model, specifically `llama2` via Ollama, with no external API dependencies. The model was served locally, ensuring reproducibility and reducing potential external factors affecting performance.
</text>
\section{Results}
\section{Results}

The experiment evaluated the sensitivity of a sentiment classifier to negation words in short movie reviews from the SST-2 dataset. The model, based on a pre-trained language model, was tested on both original and negated versions of five reviews. The results showed that the model achieved 100\% accuracy on the original texts, correctly predicting the sentiment labels for all five examples. In contrast, the accuracy on the negated texts was 60\% (3 out of 5 examples correctly classified). This suggests that the model's performance is affected by the removal of negation words, indicating a dependence on such linguistic cues for sentiment classification.

However, the small sample size limits the generalizability of these findings. A larger-scale evaluation on a more extensive subset of the SST-2 dataset would be necessary to draw robust conclusions. Additionally, the analysis did not include the computation of F1-scores or other metrics, which could provide further insight into the model's sensitivity to negation. The results also highlight the importance of ensuring that the evaluation process is valid and that the model's outputs are correctly interpreted.
\section{Discussion}
\section{Discussion}

The results presented in this study indicate that sentiment classifiers, particularly those based on pre-trained language models, are sensitive to negation words in short movie reviews from the SST-2 dataset. The model demonstrated perfect accuracy (100\%) on the original texts, which contain negation words, while its performance dropped to 60\% on negated texts, where such words were systematically removed. This suggests that the model relies heavily on negation cues to infer sentiment, which is a common phenomenon in sentiment analysis tasks where negation can significantly alter the polarity of a sentence.

However, the small sample size of only five examples limits the statistical significance of these findings. The results may not be representative of the model's behavior across the entire dataset, and further investigation with a larger and more diverse subset of the SST-2 data is necessary to draw robust conclusions. Additionally, the use of accuracy as the sole performance metric may not fully capture the model's sensitivity to negation, as metrics such as F1-score could provide a more nuanced understanding of precision and recall in classification tasks.

The observed sensitivity to negation words aligns with prior research indicating that negation is a critical linguistic feature that affects sentiment polarity. However, the extent to which models can detect and account for such negation remains a challenge, particularly in short and ambiguous texts. This study highlights the importance of incorporating negation-aware mechanisms in sentiment classification models to improve their robustness and accuracy in real-world applications.

In conclusion, while the findings suggest that the model is sensitive to negation words, further work is required to generalize these results and explore the underlying mechanisms that drive this behavior. Future research should focus on expanding the experimental scope, incorporating more comprehensive evaluation metrics, and investigating the effectiveness of negation-aware models in sentiment analysis tasks.
</text>
\end{document}